\emph{Origin and numbering.} Clauses (i) and (ii) are the lemmas (iii) and (iv) of Theorem synthesis.collapsing-hole-flux-from-ground-state-on-the-ingoing-sheet, stated there for the carrier on the ingoing generator of a Schwarzschild horizon and the Killing flow $\beta_{\tau}$, which acts on it as the dilation $\delta_{\varepsilon}$ with $\varepsilon = e^{-\kappa\tau}$. They are statements about kernels on one generator and use no premise; the arguments below are moved from that proof, with $\beta_{\tau}$ written as $\delta_{\varepsilon}$ and the limit $\tau \to \infty$ as $\varepsilon \to 0^{+}$. The null vacuum and its kernel are those of Theorem \claimref{synthesis.null-vacuum-on-a-horizon-generator}(ii), and the dilation is an automorphism by its (i).

\emph{(i)} Fix $U_{0}$ and write $U = U_{0} + x$, $U' = U_{0}$, $a = \xi'(U_{0})$, $b = \tfrac12\xi''(U_{0})$, $c = \tfrac16\xi'''(U_{0})$. Then $\xi(U) - \xi(U') = ax + bx^{2} + cx^{3} + O(x^{4})$ and $\xi'(U)\xi'(U') = a(a + 2bx + 3cx^{2}) + O(x^{3})$, so the squared ratio of the weight-$\tfrac12$ kernels, $\xi'(U)\xi'(U')\,x^{2}/(\xi(U) - \xi(U'))^{2}$, is
\[
\frac{a^{2} + 2abx + 3acx^{2}}{a^{2} + 2abx + (b^{2} + 2ac)x^{2}} + O(x^{3}) \;=\; 1 + \frac{ac - b^{2}}{a^{2}}x^{2} + O(x^{3}) \;=\; 1 + \frac{S_{\xi}(U_{0})}{6}x^{2} + O(x^{3}),
\]
because $(ac - b^{2})/a^{2} = \tfrac16\xi'''/\xi' - \tfrac14(\xi''/\xi')^{2} = \tfrac16 S_{\xi}$. The ratio itself is therefore $1 + \tfrac{1}{12}S_{\xi}x^{2} + O(x^{3})$, and the difference of the kernels is $\bigl(\tfrac{1}{12}S_{\xi}x^{2} + O(x^{3})\bigr)/x = \tfrac{1}{12}S_{\xi}(U_{0})\,x + O(x^{2})$: continuous across the diagonal and zero on it. The $i0$ prescriptions agree because $\xi' > 0$ preserves the orientation. This covariance is specific to the kernel and its weight: the same kernel $1/(U - U' + i0)$ carried with weight $0$ or $1$ has squared ratio $x^{2}/(\xi(U) - \xi(U'))^{2} = a^{-2} + O(x)$, respectively $(\xi'(U)\xi'(U'))^{2}x^{2}/(\xi(U) - \xi(U'))^{2} = a^{2} + 2abx + O(x^{2})$, neither $1 + O(x^{2})$ unless $a = 1$ and $b = 0$, so for those weights the difference is singular on the diagonal unless $\xi' \equiv 1$ (a test of the covariance weight, not of other fields). A weight-$h$ field with its own ground-state kernel $(U - U' + i0)^{-2h}$ has ratio $(1 + \tfrac16 S_{\xi}x^{2} + O(x^{3}))^{h}$, so the difference of its kernels is $\tfrac{h}{6}S_{\xi}(U_{0})x^{2 - 2h} + O(x^{3 - 2h})$, continuous for every $h \le 1$ and equal to $S_{\xi}/6$ on the diagonal at $h = 1$ (checked: the squared ratio's $x^{2}$ coefficient is $S_{\xi}/6$, which is the weight-$1$ difference's diagonal value). Where $\xi'(U_{0}) = 0$ (for example $\xi = U^{3}$ at $0$) the ratio vanishes. The expansion is exact algebra in $a, b, c$; the check computes the squared ratio as an exact rational function of $x$ for three polynomial $\xi$ at three base points, tabulating $S_{\xi}/6$, with the covariance-weight controls. Consequently, if a kernel is $W_{0}$ in the parameter $U$ plus a remainder bounded near a point, it is $W_{0}$ in the parameter $\xi$ plus a remainder bounded there (the Schwarzian term is continuous, hence bounded), and conversely.

\emph{(ii)} For $g \in L^{1}\cap L^{2}$ let $\psi(g) = \int g(U)\psi(U)\,dU$. Then $\delta_{\varepsilon}(\psi(g)) = \psi(g_{\varepsilon})$ with $g_{\varepsilon}(U) = \varepsilon^{-1/2}g(U/\varepsilon)$, the unitary dilation of $L^{2}(\R, dU)$, whose support shrinks to $U = 0$ as $\varepsilon \to 0^{+}$. Then
\[
\omega\bigl(\delta_{\varepsilon}(\psi^{*}(f)\psi(g))\bigr) \;=\; \iint \bar f_{\varepsilon}(U')g_{\varepsilon}(U)\,\bigl[W_{0}(U, U') + R(U, U')\bigr]\,dU\,dU'
\;=\; \iint \bar f(x')g(x)\,W_{0}(x, x')\,dx\,dx' \;+\; \varepsilon\iint \bar f(x')g(x)\,R(\varepsilon x, \varepsilon x')\,dx\,dx' ,
\]
substituting $U = \varepsilon x$, $U' = \varepsilon x'$: the measure gives $\varepsilon^{2}$, the two half-density factors $\varepsilon^{-1}$, and $W_{0}(\varepsilon x, \varepsilon x') = \varepsilon^{-1}W_{0}(x, x')$ (the $i0$ is unchanged for $\varepsilon > 0$), so the singular part is exactly invariant, while the remainder keeps one factor $\varepsilon$ (checked: for a polynomial remainder the scaled contribution is a polynomial in $\varepsilon$ with no constant term; for weight $0$ or $1$ the singular part is multiplied by $\varepsilon^{-1}$ or $\varepsilon$ and is not invariant). Since $R$ is bounded on a neighbourhood of $(0, 0)$ by some $C$, for $\varepsilon$ small enough that $\varepsilon\,\mathrm{supp}$ lies in that neighbourhood (test functions of compact support, dense in $L^{1}\cap L^{2}$; the general case by approximation, the two-point functions being bounded on $L^{2}$), the second term is at most $\varepsilon C\|f\|_{1}\|g\|_{1} \to 0$. So $\omega\circ\delta_{\varepsilon} \to \omega_{0}$ on every two-point function; a gauge-invariant quasi-free state is determined by its two-point function (Fact \claimref{hawking.imported-quasi-free-kms-reduces-two}, the quasi-free structure), and $\omega\circ\delta_{\varepsilon}$ is gauge-invariant quasi-free with two-point function converging to that of $\omega_{0}$, so $\omega\circ\delta_{\varepsilon} \to \omega_{0}$ on every monomial, hence weakly on the carrier. The hypothesis is sufficient and not necessary: all the proof uses is $\varepsilon\iint\bar f g R(\varepsilon x, \varepsilon x') \to 0$, which also holds when $R$ is bounded but discontinuous at the origin, and for every finite-rank change of the symbol $T_{0} \mapsto T_{0} + \sum_{k}c_{k}|a_{k}\rangle\langle b_{k}|$ with $a_{k}, b_{k} \in L^{2}$, since the dilated test functions $f_{\varepsilon}$ tend weakly to zero as $\varepsilon \to 0$ and so $\langle a_{k}, f_{\varepsilon}\rangle \to 0$; such states lie outside the class of (i) when the $a_{k}$ are discontinuous at $0$ and yet have the limit $\omega_{0}$. If instead $R$ is homogeneous of degree $-1$, $R(\varepsilon x, \varepsilon x') = \varepsilon^{-1}R(x, x')$, the second term equals $\iint\bar f g R$ for every $\varepsilon$ and does not vanish (checked with a rational homogeneous kernel); such an $R$ is unbounded near the origin.
